\documentclass[10pt]{article}

\usepackage[utf8]{inputenc}

\usepackage[T1]{fontenc}

\usepackage{amsmath}

\usepackage{amsfonts}

\usepackage{amssymb}

\usepackage[version=4]{mhchem}

\usepackage{stmaryrd}

\usepackage{graphicx}

\usepackage[export]{adjustbox}

\graphicspath{ {./images/} }

\usepackage{bbold}



\begin{document}

ранствах $(E, G)$, наделенных нек-рой топологией, согласованной с $G$, являются так наз. $G-0 \mathrm{p}$ т о н о рм и р о в а н ны е ба 3 и с ы в $E$, т. е. такие базисы $\left\{e_{n}\right\}$ топологического векторного пространства $E$, что $\left(G e_{k}, e_{n}\right)= \pm \delta_{k n} ; k, n=1,2, \ldots$ (cM. [4]).



Лum.: [1] М а ль и е в А. И., Основы линейной алгебры, Cp матем, 1944 т. 8 c. $243-80$. [3]. С., "Изв. АH CCCP. К р е й н М. Г., “Тр. Моск. матем. об-ва», 1956, т. 5, с. $367-$



\begin{center}

\includegraphics[max width=\textwidth]{2023_07_08_2a123c862ca46159a9b4g-1(1)}

\end{center}



\includegraphics[max width=\textwidth, center]{2023_07_08_2a123c862ca46159a9b4g-1}

в кн.: Вторая летняя матем. школа, ч. 1, К., 1965 , с. 15-92; [6] 1971 , т. 26 , в. 4, с. 43-92; [7] Н а д ь К., Пространства состояний с индефинитной метрикой в квантовой теории поля, пер. с англ., М., 1969; [8] И о х в и д о в И. С., “Изв. АН МолIд. ССР", 1968, № 1, с. $60-80 . \quad$ H. К. Никольскии, В. С. Павлов.



ПРОСТРАНСТВО $\mathrm{C}$ МЕРОЙ $(X, A, \mu)$ - измеримое пространство $(X, A)$ с заданной на $A$ мерой $\mu$ (т. е. счетно аддитивной функцией со значениями в $[0, \infty]$, для к-рой $\mu(\varnothing)=0 ;$ последнее свойство следует из аддитивности, если мера конечна, т. е. не принимает значения $\infty$, и даже если имеется хоть одно $Y \in A$ с $\mu Y<\infty)$. Запись $(X, A, \mu)$ часто сокращают до $(\bar{X}, \mu)$ и говорят, что $\mu$ есть мера на $X$; иногда әту запись сокращают даже до $X$. Основной случай - когда $A$ является $\sigma$-алгеброй и $X$ можно представить в виде $\bigcup_{n=1}^{\infty} X_{n}$ с $X_{n} \in A$ и $\mu X_{n}<\infty ;$ меру в этом случае наз. (в п о л н е) $\sigma-к о$ н е ч но й (а если $\mu X<\infty$, то (в по л н е) к ов е ч о й). Такова, напр., мера Лебега на $\mathbb{R}$ (см. Лебега пространство). Однако иногда встречаются и не $\sigma$-конечные меры, как, напр., $k$-мерная $X$ аусдорфа мера на $\mathbb{R}^{n}$ при $k<n$. Встречаются также модификации, когда $\mu$ принимает значения в $(-\infty$, $\infty]$, комплексные или векторные значения, а также когда $\mu$ всего лишь конечно аддитивна.



Лит.: [1] Х а л м о ш П., Теория меры, пер. с англ., М., 1953; [2] Д а н ф о р д Н., ІІ в а р ц Д ж., Јинейные оператары, ч. 1 - Общая теория, пер. с англ., М., 1962. Д. В. Аносов.



ПРОСТРАНСТВО-ВРЕМЯ - термин, обозначающий геометрич. конструкцию, к-рая описывает пространственные и временны́е отношения в тех физич. теориях, в к-рых эти отношения рассматриваются как взаимозависящие (эти теории принято наз. релятивистскими). Впервые понятие П.-в. возникло при формулировке и систематизации основных положений теории относительности. П.-в. этой теории является qетырехмерным псевдоевклидовым пространством $E_{(1,3)}^{4}$ с линейным элементом



\[

d s^{2}=c^{2} d t^{2}-d x^{2}-d y^{2}-d z^{2}

\]



где $x, y, z$ - пространственные координаты, а $t$ временна́я координата, $c$ - скорость света. Эта система координат называется в физике галилеевой системой координат и соответствует инерциальной системе отсчета. Переход между различными галилеевыми системами координат, соответствуюгий рассмотрению инерциальных систем отстета, движущихся друг относительно друга, осуществляется с помощью Лоренца преобразований. То, что значение временно́й координаты в новой системе координат оказывается при этом выраженным как через временну́ю, так и через пространственные координаты старой системы координат, отражает взаимозависимость простравственных и временны́х отношений в специальной теории относительности. П.-в. специальной теории относительности принято называть также П.-в. М и н к о в с к о го, ло р е нце в ы м П.-B.



В общей теории относительности в качестве П.-в. используются различные четырехмерные псевдоримановы пространства сигнатуры $(1,3)$. Отличие метрики этого П.-в. от плоской метрики П.-в. сшециальной теории относительности описывает в общей теории относительности гравитационное поле (см. Гравитация). В свою очередь метрика П.-в. связана с распределени- ем и свойствами негравитационных полей и различных видов с помощью Эйнштейна уравнений.



Выработка концепции П.-в. сыграла важную роль в преодолении метафизич. подхода к пространству как абсолютному вместилищу тел и к времени как абсолютной длительности, не связанной с реальвыми физич. процессами.



В дальнейшем концепция П.-в. в той или иной форме входит в структуру других физич. теорий, рассматривающих релятивистские әффекты (релятивистской квантовой механики, квантовой теории поля и др.). В общей теории относительности были изучены многие типы П.-в., являющиеся решениями уравнений Эйнштейна.



Качественное различие между пространственными и временны́ми отношениями с точки зрения релятивистской физики находит свое отражение в наличии в П.-в. векторов различной природы - пространственнои времениподобных векторов, образующих в касательных пространствах конусы. Соответственно, метрика П.-в. оказывается индефинитной, пространственно- и времениподобные векторы имеют разные знаки скалярного квадрата. Границу между конусами пространственно- и времениподобных векторов образует изотропный конус, изотропные векторы к-рого имеют нулевой скалярный квадрат и соответствуют движению света и других частиц с нулевой массой покоя.



Многие специфич. әффекты теории относительности связаны именно с индефинитностью метрики П.-в. И с наличием структуры изотропных конусов на П.-в. Напр., лоренцево замедление времени есть следствие обратного неравенства треугольника в пространстве с индефинитной метрикой, согласно к-рому в двумерном псевдоевклидовом пространстве наклонная всегда короче своей проекции.



Во многих случаях оказывается полезным в различной степени отвлекаться от конкретного строения метрики П.-в. и рассматривать лишь свойства структуры изотропных конусов на П.-в., то есть рассматривать различные т. н. общие пространства кинематич. типа, или времениподобные пространства.



При ретроспективном анализе предшествующих физич. теорий с точки зрения теории относительности были построены различные типы П.-в., к-рые можно условно поставить в соответствие ньютоновской механике (галилеево пространство) и даже физич. представлениям Аристотеля (см. [5]). Эти П.-в. являются различными пространствами с вырожденным изотропным (световым) конусом (напр., полуримановым пространством). Именно вырождение изотропного конуса позволяет рассматривать әти пространства как предельные случаи П.-в. теории относительности и сопоставлять их тем теориям, исходной внутренней структуре к-рых концепция П.-в. была чужда.



Лит.: [1] Э й н ш т е й н А., Собр. науч. трудов, т. 1-4, М , 1965-67; [2] А л е к с а н д р о в А. Д., "Вопросы философии", $1959, ﹎{0} 1$, с. $67-84$; [3] Л а н д а у Ј. Д., ЈІ и фш и ц Е. М., Теория поля, 6 изд., М., 1973 (Теоретич. физика, т. 2); [4] Р а ш е в с к й й П. К., Риманова геометрия и тензор-

ный анализ, 3 изд., М., 1967; [5] П е н р о у з Р., Структура проный анализ, 3 изд., М., 1967; [5] П е н р о у з Р., Структура про-

странства-времени, пер. с англ., М., 1972. Д. Соколов.



ПРОТИВОПОЛОЖНАЯ ТЕОРЕМА - теорема, получающаяся путем замены условия и заключения данной исходной теоремы их отрицаниями.



ПРОТИВОРЕЧИВЫЙ КЛАСС - класс $K$ формул языка узкого исчисления предикатов (УИП) такой, что существует такая формула $\varphi$, что средствами УИП из $K$ выводимо как $\varphi$, так и $7 \varphi$ (отрицание $\varphi$ ). Другими словами, если к аксиомам УИП добавить все формулы из $K$ в качестве новых аксиом, то в полученном исчислении будет выводима как формула $\varphi$, так и формула $7 \varphi$.



B. Н. Гричичн. ПРОТИВОРЕЧИЕ - формула $\varphi$ языка узкого исчисления предикатов (УИП) такая, что во всех моделях





\end{document}